% C05 replacement for the acquisition subsection's equation and explanation.
% Integrate before the unchanged scalarized-baseline paragraph.
Let $\boldsymbol{\xi}$ be a joint posterior draw at the observed baseline and
at the proposed batch $\mathbf{X}$. Define the feasible sampled baseline
$\mathcal{P}_{\mathrm{feas}}(\boldsymbol{\xi}_n)$ and the feasible candidate
subset $\mathcal{F}(\boldsymbol{\xi},\mathbf{X})$ by requiring all eight
inequalities at each individual point. The ideal constrained noisy
hypervolume-improvement criterion is
\begin{equation}
 \alpha_{\mathrm{ideal}}(\mathbf{X})=
 \log\mathbb{E}_{\boldsymbol{\xi}}\!\left[
 \mathrm{HVI}\!\left(\mathcal{F}(\boldsymbol{\xi},\mathbf{X})\mid
 \mathcal{P}_{\mathrm{feas}}(\boldsymbol{\xi}_n),\mathbf{r}_{\mathrm{acq}}\right)
 \right].
 \label{eq:qlognehvi}
\end{equation}
This expression defines the ideal target; the implementation uses BoTorch's
qLogNEHVI approximation~\cite{Daulton2021,Ament2023}. It evaluates batch
hypervolume contributions by inclusion--exclusion and applies feasibility
weights to each candidate participating in a contribution, rather than
multiplying the entire batch improvement by one feasibility factor. Sampled
baseline feasibility determines the baseline front. For each sampled candidate,
the weight is $\prod_{k=1}^{8}h(-c_k/\tau_c)$, where $h$ is BoTorch's
fat-tailed sigmoid and $\tau_c=10^{-3}$. Positive improvements and extrema are
also smoothed and combined in log space, with $\tau_{\mathrm{ReLU}}=10^{-6}$
and $\tau_{\max}=10^{-3}$. The unconstrained variant omits the diagnostic
constraints. The separately available analytical Gaussian probability of
feasibility is a diagnostic; it is not substituted for these sampled weights.

The residuals are expressed in their recorded physical units: metres, radians,
electron-volts and transmission fraction, without threshold normalization.
Thus $\tau_c$ corresponds to \SI{1}{\milli\metre},
\SI{1}{\milli\radian}, \SI{0.001}{\electronvolt} and a transmission
fraction of 0.001 in their respective channels. These are scale parameters,
not equal physical transition widths. At a residual of zero the indicator is
0.5; at residuals $-\tau_c$ and $+\tau_c$ it is approximately 0.8089 and
0.1911, respectively. For a \SI{1}{\milli\metre} beam-size limit,
a size of zero therefore receives an indicator of 0.8089, while a size of
\SI{2}{\milli\metre} receives 0.1911. This unit dependence limits the
interpretation of constraint weighting; no residual scaling is changed when
analysing the recorded campaigns.

The executed call sites use baseline pruning and BoTorch's default normal
Sobol samplers: 128 posterior samples for qLogNEHVI and 512 for scalarized
qLogNEI. These counts and the smoothing constants are reconstructed from the
call sites and the matching archived BoTorch version, because acquisition
objects and sampler base draws were not saved. The model dimension and baseline
sizes remain below the Sobol engine's dimension limit. Scalarized qLogNEI uses
$\tau_{\max}=10^{-2}$ and $\tau_{\mathrm{ReLU}}=10^{-6}$. These posterior
sample counts differ from the 1,024 raw candidate-initialization samples.
